\documentclass{article}
\usepackage{amsmath,amssymb,amsthm}
\begin{document}

\textbf{Subproblem 4 --- Small-case classification: triangles fail, and the only
quadrilateral extremizer is a parallelogram.}

\medskip
\textbf{Context.}
Use the full setup: $S\subset\mathbb{R}^2$ finite, $|S|\ge 3$, $S$ in convex
position, $T$ a finite set with both signs, $\varepsilon:T\to\{+1,-1\}$,
and define $F$ as before.  Write $P=\operatorname{conv}(S)$,
$Q=\operatorname{conv}(T)$.

By Subproblem~1 (which may be freely used), every vertex of $P+Q$ lies in
$\operatorname{supp}(F)$, so $|\operatorname{supp}(F)|\ge|\operatorname{vert}(P+Q)|$.

By Subproblem~3, if $|\operatorname{supp}(F)|=n$ then $n\le 4$ (may also be freely
used).  Together with $n\ge 3$, we have $n\in\{3,4\}$.

\medskip
\textbf{Statement.}  Prove:
\begin{enumerate}
  \item \emph{(Triangle case impossible.)}  If $n=3$ (so $P$ is a triangle) and
        $|T|\ge 2$, then $|\operatorname{vert}(P+Q)|\ge 4 > 3 = n$, so
        $|\operatorname{supp}(F)|\ge 4$ and equality $|\operatorname{supp}(F)|=3$ is
        impossible.
  \item \emph{(Quadrilateral must be a parallelogram.)}  If $n=4$ (so $P$ is a
        convex quadrilateral) and $|\operatorname{supp}(F)|=4$, then $S$ is the
        vertex set of a parallelogram.
\end{enumerate}

\medskip
\textbf{Hints for (1).}
$P$ has exactly $3$ outer edge-normal directions.  $Q=\operatorname{conv}(T)$ is
a non-degenerate convex set (since $T$ has at least $2$ points with $|T|\ge 2$).
If $Q$ is a segment, its one edge-normal direction is either among the $3$
normals of $P$ or is new.
\begin{itemize}
  \item \emph{New direction:} $P+Q$ has $3+2=5$ vertices $>3$.
  \item \emph{Direction parallel to edge $e$ of $P$:} The Minkowski sum
        $P+Q$ ``stretches'' edge $e$ but also creates \emph{two new vertices} at
        the ends of the stretched edge (the two endpoints of $Q$ each combine
        with the two endpoints of $e$, but only the outer two survive as new
        vertices of the hull). Concretely, if $e=v_1v_2$ and $Q=[0,w]$ with $w$
        parallel to $e$, the hull of $P+Q$ acquires the edge $v_1 v_1+w$ (or
        $v_2 v_2+w$), giving $4$ vertices.
\end{itemize}
In all cases $|\operatorname{vert}(P+Q)|\ge 4$.
If $Q$ is $2$-dimensional (a non-degenerate polygon), use the fact that a
$2$-dimensional addend always contributes new edge-normal directions beyond
those already in $P$ (a triangle), giving at least $4$ vertices.

\medskip
\textbf{Hints for (2).}
Use the rigidity from Subproblem~1: $|\operatorname{supp}(F)|=4$ implies
$|\operatorname{vert}(P+Q)|=4$.
The number of vertices of the Minkowski sum of two convex polygons equals the
number of distinct outer-normal directions in their combined normal fans.
$P$ is a quadrilateral with $4$ outer-normal directions.  For $P+Q$ to have
exactly $4$ vertices, $Q$ must contribute \emph{no new} outer-normal directions:
every edge direction of $Q$ is parallel to some edge of $P$.

In particular, $Q$ can be at most a \emph{centrally symmetric hexagon} (or
degenerate versions thereof) whose edges are parallel to the edges of $P$.
But examine what it means for the surviving lonely points to be exactly
$4$: the four vertices of $P+Q$ must come from the four vertices of $P$
each paired with a corresponding vertex of $Q$.  This forces the pairs to
satisfy $v_1+q_1,\, v_2+q_2,\, v_3+q_3,\, v_4+q_4$ (in cyclic order), and the
transition pattern between exposed $Q$-vertices around the boundary of
$P+Q$ forces $Q$ to alternate between two vertices only (i.e.\ $Q$ is a
segment).

With $Q=[0,v]$ a segment parallel to edge $v_1v_2$ of $P$ (and also parallel to
the opposite edge $v_3v_4$), the four surviving vertices of $P+Q$ are
$v_1,\,v_2+v,\,v_3+v,\,v_4$ (or similar).  But for $P+Q$ to have only $4$
vertices, the edge $v_2v_3$ of $P$ must be parallel to the edge $v_4v_1$
(otherwise those edges' normal directions are new and create more vertices).
Parallel opposite sides means $P$ is a parallelogram.

\end{document}
